\hypertarget{introduction}{%
\chapter{Introduction}\label{introduction}}

This chapter provides an overview of the IDRISK2 project, including its motivation, objectives, and the overarching problem it aims to solve within the context of European food safety regulations and data management.

\hypertarget{context-and-problem-statement}{%
\section{Context and Problem Statement}\label{context-and-problem-statement}}

Food security represents a critical challenge, with significant implications for public health, economic stability, and societal well-being, particularly in the context of increasingly complex and globalised food supply chains. This complexity is further exacerbated by emerging food safety hazards stemming from factors such as climate change, novel materials in agri-food systems, and evolving dietary trends, necessitating advanced systems for their identification, assessment, and management \parencite{john_f__leslie_d22f2c65}.

For this reason, the European Food Safety Authority (EFSA) has developed standardised frameworks such as FoodEx2 to classify and describe food products consistently, which is fundamental for effective food safety risk assessment \parencite{hern_n_g__redondo_919f3cec}. FoodEx2 is a hierarchical system that represents food items using base terms and facets, which add further information about the food items.

Within the FoodEx2 hierarchy, three main types of terms are defined: hierarchy terms (aggregated groups), core terms (representing the minimum level of detail for coding), and extended terms (providing additional detail beyond core terms) ((EFSA), 2015). In practice, food coding primarily relies on core and extended terms, collectively referred to as base terms. To further enhance the level of specificity, FoodEx2 incorporates facets, which describe additional characteristics of food items, such as origin, processing methods, physical state, and other relevant attributes. Additionally, some terms include implicit facets that are not explicitly codified but are inherently understood from the base term itself, such as those related to the source facet and the part-nature facet ((EFSA), 2015).

The FoodEx2 represents a significant advancement in food classification and monitoring, enabling the standardised representation of food items through a harmonised coding scheme. Nevertheless, the manual mapping of food descriptions to FoodEx2 codes remains a time-consuming and error-prone process, often leading to inconsistencies in data reporting. This laborious manual classification not only hinders efficient data analysis but also complicates the comparison of food consumption data across diverse sources, ultimately impacting the accuracy of food safety risk predictions (Eftimov et al., 2017).

The inherent complexity of manually classifying vast quantities of diverse food data according to the intricate FoodEx2 hierarchy often leads to inconsistencies, delayed regulatory responses, and significant operational costs for food safety authorities and industry stakeholders. In this context, recent advances in artificial intelligence have spurred the development of innovative solutions aimed at improving the efficiency and accuracy of food classification. Nevertheless, a persistent methodological debate remains regarding the optimal approach for hierarchical classification within multi-dimensional taxonomies, specifically, the choice between fine-tuning versus retrieval-augmented strategies for mapping natural language inputs \parencite{lorenzo_molfetta_f41ce012}.

Moreover, despite this progress, AI models continue to encounter significant limitations, particularly when interpreting complex multimodal data while simultaneously ensuring adherence to strict data privacy regulations, such as the General Data Protection Regulation. These challenges underscore the necessity for further research into developing AI-enhanced systems tailored specifically for FoodEx2.

Beyond these challenges, advanced AI models are susceptible to hallucinations, a phenomenon characterised by the generation of plausible yet factually incorrect or fabricated information, particularly in tasks that require accurate factual knowledge and contextual understanding. This vulnerability is particularly concerning in regulatory environments where the precision of automated risk assessment directly impacts public health outcomes \parencite{omer_faruk_sari_c479eead}. Furthermore, these models often require extensive training periods, which may limit their applicability in dynamic environments where continuous updates are necessary. This limitation is particularly relevant in the food domain, as the FoodEx2 classification system undergoes periodic updates, thus requiring adaptable and robust AI architectures.

\hypertarget{motivation}{%
\section{Motivation}\label{motivation}}

Motivated by the challenges in food classification, the IDRISK2 project addresses a critical bottleneck by proposing an AI-enhanced digital registration and analysis system designed to automate and streamline the process of FoodEx2 compliance. This system aims to improve data quality and facilitate more accurate risk assessments. Furthermore, it seeks to promote greater harmonisation of data collected by EFSA, contributing to more reliable and consistent food safety analyses.

Beyond its regulatory impact, this project also addresses operational challenges, where the automation of food classification can significantly reduce the time and manual effort required from inspectors, increasing efficiency and reducing inconsistencies. In this context, IDRISK2 aims to optimise the workflow of the registration and analysis process, enabling resources to be allocated more strategically and proactively in risk detection activities.

Additionally, the project contributes to technological innovation by exploring new use cases of artificial intelligence in food safety, combining Optical Character Recognition (OCR), Computer Vision, Vision-Language Models (VLMs), and Retrieval-Augmented Generation (RAG). The integration of these technologies enables digital registration systems to interpret and classify complex multimodal data more autonomously and accurately, improving the reliability of food classification processes.

Finally, this research aims to provide scientific contributions by advancing the understanding of the applicability of multimodal AI models and RAG approaches in complex regulatory environments. This includes evaluating their capabilities and limitations in interpreting documents and images for compliance with standards such as FoodEx2. Furthermore, this study adopts a more nuanced perspective on the regulatory challenges posed by the EU AI Act, recognizing that compliance requires going beyond simple log registration to address complex issues of risk classification and transparency \parencite{idse_lucien_val_e2ea81c6}. By fostering critical discussion on these requirements, this research contributes to the development of future, more robust approaches in food safety systems.

\hypertarget{research-questions}{%
\section{Research Questions}\label{research-questions}}

Based on the challenges and motivations identified in the previous sections, the following central research question is proposed.

How can the integration of Artificial Intelligence techniques, combining computer vision, Optical Character Recognition, and large language models enhanced with Retrieval-Augmented Generation, improve the efficiency, accuracy, and security of the FoodEx2 coding process in digital food sample registration, ensuring compliance with EFSA standards and relevant regulatory cybersecurity requirements for food safety systems?

\hypertarget{research-objectives}{%
\section{Research Objectives}\label{research-objectives}}

To address the central research question, this dissertation aims to develop and evaluate an AI-driven system for automated FoodEx2 classification, ensuring compliance with EFSA standards and regulatory requirements.

To achieve this goal, the following specific objectives are defined:

\begin{itemize}
\item
  Design a multimodal architecture integrating Optical Character Recognition (OCR), Computer Vision, and Large Language Models for food data interpretation
\item
  Develop a Retrieval-Augmented Generation (RAG) pipeline to support reliable and explainable FoodEx2 classification
\item
  Ensure compliance with GDPR and cybersecurity requirements for secure data processing
\item
  Evaluate the system performance using classification metrics such as accuracy, precision, recall, and processing time
\item
  Compare the proposed approach with baseline methods, including manual classification and non-RAG AI approaches
\item
  Assess the system's scalability and adaptability to evolving FoodEx2 standards
\end{itemize}

\hypertarget{scope-and-limitations}{%
\section{Scope and Limitations}\label{scope-and-limitations}}

Based on the motivations and objectives outlined, the scope of this research is limited to the development and evaluation of an AI-enhanced system tailored for FoodEx2 compliance within the context of digital food sample registration. This includes the classification of food items and ingredients according to the hierarchical structure of FoodEx2, incorporating both base terms and facets, while prioritizing the secure handling of multimodal data.

Additionally, this research includes the development of a digital registration workflow that supports data capture at the point of sampling and the generation of structured FoodEx2 outputs. The evaluation will consider classification performance and usability within simulated or controlled scenarios aligned with regulatory workflows.

The proposed system focuses on the use of artificial intelligence techniques to automate the interpretation of textual and visual information associated with food samples. This includes the extraction of textual information from product labels using Optical Character Recognition and the analysis of visual characteristics through Computer Vision techniques. Furthermore, the system explores the integration of Large Language Models with a Retrieval-Augmented Generation framework to enhance classification accuracy, explainability, and regulatory compliance.

Moreover, the scope includes the development of a secure, closed AI model environment, ensuring data privacy and regulatory compliance by avoiding reliance on public cloud solutions or external unverified APIs. The system architecture will prioritize adherence to GDPR and other relevant cybersecurity protocols, ensuring robust data protection throughout the registration and analysis pipeline. However, it is important to note that, as a research prototype, this system focuses on foundational security and privacy principles rather than claiming full certification or compliance with the complex "high-risk" obligations stipulated by the EU AI Act (cf. Val, 2025 \parencite{idse_lucien_val_e2ea81c6}).

However, this research acknowledges certain limitations. The inherent complexity and continuous evolution of the FoodEx2 taxonomy may require ongoing model retraining and adaptation, which fall outside the scope of this dissertation. Additionally, the proposed system will be developed as a research prototype and may not include full integration with existing regulatory infrastructures such as EFSA or LIMS systems. Furthermore, the evaluation will be conducted using selected datasets and scenarios, which may not fully represent the diversity of real-world food products and operational conditions.
